% ============================================================
\chapter{Cycles Limites et Th\'eor\`eme de Poincar\'e-Bendixson}
\label{ch:cycles_limites}
% ============================================================

Le c\oe ur bat. Un circuit \'electronique oscille. Une population de pr\'edateurs et de proies fluctue p\'eriodiquement. Ces ph\'enom\`enes d'oscillation auto-entretenue sont omnipr\'esents dans la nature et la technologie. Math\'ematiquement, ils correspondent \`a des \emph{cycles limites}~: des orbites p\'eriodiques isol\'ees vers lesquelles les trajectoires voisines convergent (ou divergent). Le th\'eor\`eme de Poincar\'e-Bendixson (1901) est le r\'esultat fondamental qui garantit l'existence de tels cycles dans les syst\`emes plans~: si une trajectoire reste confin\'ee dans une r\'egion compacte sans point d'\'equilibre, elle doit converger vers un cycle limite. Ce r\'esultat, sp\'ecifique \`a la dimension 2, n'a pas d'analogue en dimension sup\'erieure --- c'est l\`a que le chaos entre en sc\`ene.

\section{Introduction}

Les cycles limites sont des orbites p\'eriodiques isol\'ees dans l'espace des
phases, c'est-\`a-dire des orbites ferm\'ees qui ne font pas partie d'une
famille continue d'orbites ferm\'ees. Ils constituent l'un des ph\'enom\`enes
les plus importants propres aux syst\`emes non lin\'eaires, car les syst\`emes
lin\'eaires ne peuvent produire que des centres (familles continues d'orbites
ferm\'ees). Le th\'eor\`eme de Poincar\'e-Bendixson fournit un crit\`ere
puissant pour prouver l'existence de cycles limites dans le plan.

\section{D\'efinitions}

\begin{definition}[Cycle limite]
Un \textbf{cycle limite} est une orbite p\'eriodique $\Gamma$ qui est
l'ensemble $\omega$-limite ou $\alpha$-limite d'au moins une orbite
n'appartenant pas \`a $\Gamma$.
\end{definition}

\begin{definition}[Types de cycles limites]
Un cycle limite $\Gamma$ est :
\begin{itemize}
    \item \textbf{stable} (attractif) si $\Gamma = \omega(x_0)$ pour tout
    $x_0$ dans un voisinage de $\Gamma$ ;
    \item \textbf{instable} (r\'epulsif) si $\Gamma = \alpha(x_0)$ pour tout
    $x_0$ dans un voisinage de $\Gamma$ ;
    \item \textbf{semi-stable} si $\Gamma$ est attractif d'un c\^ot\'e et
    r\'epulsif de l'autre.
\end{itemize}
\end{definition}

\begin{exemple}[Cycle limite dans l'oscillateur de Van der Pol]
L'\'equation de Van der Pol
\[
    \ddot{x} - \mu(1 - x^2)\dot{x} + x = 0, \qquad \mu > 0,
\]
poss\`ede un unique cycle limite stable. Pour $\mu$ petit, ce cycle est
proche du cercle de rayon 2 centr\'e \`a l'origine. Pour $\mu$ grand,
le cycle prend une forme caract\'eristique d'\textbf{oscillation de relaxation}.
\end{exemple}

\section{Crit\`ere de non-existence : Bendixson-Dulac}

\begin{theoreme}[Crit\`ere de Bendixson]
\label{thm:bendixson}
Soit $D \subseteq \R^2$ un domaine simplement connexe. Si
$\mathrm{div}\, f(x) = \frac{\partial f_1}{\partial x_1} + \frac{\partial f_2}{\partial x_2}$
est de signe constant (et non identiquement nulle) sur $D$, alors le
syst\`eme $\dot{x} = f(x)$ n'a \textbf{aucune orbite p\'eriodique}
enti\`erement contenue dans $D$.
\end{theoreme}

\begin{proof}
Par le th\'eor\`eme de Green, si $\Gamma$ \'etait une orbite p\'eriodique
enfermant un domaine $\Omega$, on aurait
\[
    \iint_\Omega \mathrm{div}\, f\, dx_1\, dx_2 =
    \oint_\Gamma (f_1\, dx_2 - f_2\, dx_1) = 0,
\]
car sur $\Gamma$, $dx_i = f_i\, dt$, et l'int\'egrale curviligne se simplifie.
Mais si $\mathrm{div}\, f$ est de signe constant et non nulle, l'int\'egrale
double est non nulle : contradiction.
\end{proof}

\begin{theoreme}[Crit\`ere de Dulac]
Le m\^eme r\'esultat vaut si l'on remplace $\mathrm{div}\, f$ par
$\mathrm{div}(\rho f)$ o\`u $\rho : D \to \R$ est une fonction de classe
$C^1$ strictement positive (fonction de Dulac).
\end{theoreme}

\begin{exemple}[Syst\`eme de Lotka-Volterra]
Pour $\dot{x} = x(\alpha - \beta y)$, $\dot{y} = y(\delta x - \gamma)$ dans
le premier quadrant, prenons $\rho(x,y) = \frac{1}{xy}$. Alors
$\mathrm{div}(\rho f) = 0$, ce qui ne permet pas de conclure. Cependant,
l'existence d'une premi\`ere int\'egrale montre que les orbites sont
ferm\'ees (centre) et non des cycles limites.
\end{exemple}

\section{Le th\'eor\`eme de Poincar\'e-Bendixson}

\begin{theoreme}[Poincar\'e-Bendixson]
\label{thm:poincare-bendixson}
Soit $\dot{x} = f(x)$ un syst\`eme de classe $C^1$ sur un ouvert de $\R^2$.
Soit $\gamma^+(x_0)$ une semi-orbite positive born\'ee contenue dans un
compact $K$ ne contenant aucun point d'\'equilibre. Alors
$\omega(x_0)$ est une orbite p\'eriodique (et donc un cycle limite si
$x_0 \notin \omega(x_0)$).
\end{theoreme}

\begin{remarque}
Plus g\'en\'eralement, si $K$ contient un nombre fini de points d'\'equilibre,
l'ensemble $\omega(x_0)$ est soit :
\begin{enumerate}
    \item un point d'\'equilibre ;
    \item une orbite p\'eriodique ;
    \item un ensemble form\'e de points d'\'equilibre reli\'es par des
    connexions h\'et\'eroclines ou homoclines.
\end{enumerate}
\end{remarque}

\begin{attention}[title=\textbf{Dimension 2 uniquement}]
Le th\'eor\`eme de Poincar\'e-Bendixson est sp\'ecifique \`a la dimension 2.
En dimension $n \geq 3$, les ensembles $\omega$-limites born\'es peuvent \^etre
beaucoup plus compliqu\'es (attracteurs \'etranges, par exemple).
\end{attention}

\begin{corollaire}[Existence de cycles limites]
Soit $K \subset \R^2$ un compact positivement invariant, annulaire (sans
trou), ne contenant aucun point d'\'equilibre. Alors $K$ contient au moins
un cycle limite.
\end{corollaire}

\begin{intuition}
Strat\'egie pour prouver l'existence d'un cycle limite :
\begin{enumerate}
    \item Construire une r\'egion annulaire $K$ positivement invariante
    (le champ pointe vers l'int\'erieur sur le bord).
    \item V\'erifier qu'il n'y a pas de point d'\'equilibre dans $K$.
    \item Conclure par Poincar\'e-Bendixson.
\end{enumerate}
\end{intuition}

\section{Application de Poincar\'e et multiplicateurs}

\begin{definition}[Section de Poincar\'e]
Une \textbf{section de Poincar\'e} est une courbe $\Sigma$ transverse au
champ de vecteurs, c'est-\`a-dire telle que $f(x) \cdot n(x) \neq 0$ pour
tout $x \in \Sigma$, o\`u $n$ est un vecteur normal \`a $\Sigma$.
\end{definition}

\begin{definition}[Application de Poincar\'e]
L'\textbf{application de Poincar\'e} (ou application de premier retour)
$P : \Sigma \to \Sigma$ associe \`a un point $x \in \Sigma$ le prochain point
d'intersection de l'orbite avec $\Sigma$.
\end{definition}

\begin{proposition}
Une orbite p\'eriodique $\Gamma$ correspond \`a un point fixe de $P$.
La stabilit\'e de $\Gamma$ est d\'etermin\'ee par $\abs{P'(x^*)}$ :
\begin{itemize}
    \item $\abs{P'(x^*)} < 1$ : cycle stable ;
    \item $\abs{P'(x^*)} > 1$ : cycle instable ;
    \item $\abs{P'(x^*)} = 1$ : cas ind\'etermin\'e.
\end{itemize}
\end{proposition}

\begin{definition}[Multiplicateur de Floquet]
Pour une orbite p\'eriodique de p\'eriode $T$ du syst\`eme $\dot{x} = f(x)$
en dimension $n$, les \textbf{multiplicateurs de Floquet} sont les valeurs
propres de la matrice de monodromie
\[
    M = \Phi(T),
\]
o\`u $\Phi(t)$ est la solution de $\dot{\Phi} = Df(\gamma(t))\,\Phi$,
$\Phi(0) = I$. Un multiplicateur vaut toujours 1 (direction tangente).
L'orbite est stable si tous les autres multiplicateurs ont un module $< 1$.
\end{definition}

\section{Th\'eorie de l'indice}

\begin{definition}[Indice d'une courbe ferm\'ee]
Soit $\gamma$ une courbe ferm\'ee simple dans $\R^2$ ne passant par aucun
\'equilibre. L'\textbf{indice de $\gamma$} par rapport au champ $f$ est le
nombre de tours que fait le vecteur $f(x)$ lorsque $x$ parcourt $\gamma$
dans le sens positif :
\[
    I_\gamma = \frac{1}{2\pi}\oint_\gamma d\theta, \qquad
    \theta = \arctan\frac{f_2(x)}{f_1(x)}.
\]
\end{definition}

\begin{theoreme}[Propri\'et\'es de l'indice]
\begin{enumerate}
    \item Si $\gamma$ n'entoure aucun \'equilibre, $I_\gamma = 0$.
    \item L'indice d'un n\oe ud ou d'un foyer est $+1$.
    \item L'indice d'un col est $-1$.
    \item L'indice d'une orbite p\'eriodique est $+1$.
    \item L'indice est additif : si $\gamma$ entoure plusieurs \'equilibres,
    $I_\gamma = \sum I_{x_i^*}$.
\end{enumerate}
\end{theoreme}

\begin{corollaire}
\`A l'int\'erieur d'un cycle limite, la somme des indices des \'equilibres
vaut $+1$. En particulier, un cycle limite doit entourer au moins un
\'equilibre.
\end{corollaire}

\section{Exemples d\'etaill\'es}

\begin{exemple}[Oscillateur de Van der Pol --- existence du cycle]
Consid\'erons le syst\`eme en coordonn\'ees de Li\'enard :
\[
    \dot{x} = y - F(x), \qquad \dot{y} = -x,
\]
avec $F(x) = \mu(x^3/3 - x)$, $\mu > 0$. L'unique point d'\'equilibre est
l'origine (foyer instable pour $\mu > 0$). On construit une r\'egion
annulaire :
\begin{itemize}
    \item Bord int\'erieur : cercle de petit rayon $r$ (le champ pointe vers
    l'ext\'erieur car l'\'equilibre est instable) ;
    \item Bord ext\'erieur : courbe de niveau de l'\'energie choisie suffisamment
    grande (le syst\`eme est dissipatif globalement).
\end{itemize}
Par Poincar\'e-Bendixson, il existe un cycle limite dans cette r\'egion.
\end{exemple}

\begin{codeblock}[title=\textbf{Cycle limite de Van der Pol}]
\begin{minted}{python}
import numpy as np
from scipy.integrate import solve_ivp
import matplotlib.pyplot as plt

def van_der_pol(t, state, mu=1.0):
    x, y = state
    return [y, mu*(1 - x**2)*y - x]

fig, axes = plt.subplots(1, 3, figsize=(15, 5))
for ax, mu in zip(axes, [0.5, 1.0, 5.0]):
    for r in [0.1, 0.5, 1.0, 3.0, 5.0]:
        for theta in np.linspace(0, 2*np.pi, 4, endpoint=False):
            x0 = [r*np.cos(theta), r*np.sin(theta)]
            sol = solve_ivp(van_der_pol, [0, 50], x0,
                            args=(mu,), max_step=0.05)
            ax.plot(sol.y[0], sol.y[1], 'b-', lw=0.3, alpha=0.5)
    ax.set_title(f"$\\mu = {mu}$")
    ax.set_xlabel("$x$"); ax.set_ylabel("$\\dot{x}$")
    ax.set_xlim(-6, 6); ax.set_ylim(-8, 8)
    ax.set_aspect('auto'); ax.grid(True, alpha=0.3)
plt.tight_layout()
plt.savefig("van_der_pol_cycles.pdf")
plt.show()
\end{minted}
\end{codeblock}

\begin{exemple}[Mod\`ele de FitzHugh-Nagumo]
Le mod\`ele simplifi\'e de neurone de FitzHugh-Nagumo est
\[
    \dot{v} = v - \frac{v^3}{3} - w + I, \qquad
    \dot{w} = \epsilon(v + a - bw),
\]
avec $\epsilon \ll 1$. Ce syst\`eme poss\`ede un cycle limite pour certaines
valeurs du courant inject\'e $I$, mod\'elisant le potentiel d'action.
\end{exemple}

\section{Le th\'eor\`eme de Li\'enard}

\begin{theoreme}[Li\'enard]
Consid\'erons l'\'equation de Li\'enard $\ddot{x} + f(x)\dot{x} + g(x) = 0$
avec $F(x) = \int_0^x f(s)\, ds$. Sous les conditions :
\begin{enumerate}
    \item $f$ et $g$ continues, $g$ impaire avec $xg(x) > 0$ pour $x \neq 0$ ;
    \item $F$ impaire, $F(x) \to +\infty$ quand $x \to +\infty$ ;
    \item $F$ est n\'egative pour $0 < x < a$, positive et croissante pour
    $x > a$ (pour un certain $a > 0$) ;
\end{enumerate}
alors le syst\`eme admet un \textbf{unique cycle limite stable}.
\end{theoreme}

\section{Exercices}

\begin{exercice}[Bendixson]
Montrer que le syst\`eme $\dot{x} = x + e^{-y}$, $\dot{y} = y + e^x$ n'a
pas d'orbite p\'eriodique dans $\R^2$.
\end{exercice}

\begin{exercice}[Poincar\'e-Bendixson]
Montrer l'existence d'un cycle limite pour le syst\`eme
$\dot{r} = r(1-r^2) + \epsilon r\cos\theta$,
$\dot{\theta} = 1$ pour $\epsilon$ suffisamment petit.
\end{exercice}

\begin{exercice}[Indice]
Calculer l'indice de l'origine pour le syst\`eme
$\dot{x} = x^2 - y^2$, $\dot{y} = 2xy$.
\end{exercice}

\begin{exercice}[Application de Poincar\'e]
Pour le syst\`eme en coordonn\'ees polaires
$\dot{r} = r(1-r)(2-r)$, $\dot{\theta} = 1$, trouver tous les cycles
limites et d\'eterminer leur stabilit\'e.
\end{exercice}

\begin{exercice}[Van der Pol]
Montrer par la m\'ethode de la moyenne que pour $\mu \ll 1$, le cycle limite
de l'oscillateur de Van der Pol est proche du cercle de rayon 2.
\end{exercice}

\begin{exercice}[Unicit\'e]
Soit $\dot{x} = y$, $\dot{y} = -x + y(1 - x^2 - 3y^2)$. Montrer qu'il
existe un unique cycle limite en utilisant le th\'eor\`eme de Li\'enard.
\end{exercice}
